\documentclass[11pt]{amsart}

\usepackage{amsmath,amssymb,amsthm,mathtools}
\usepackage[margin=1.1in]{geometry}
\usepackage{enumitem}
\usepackage{hyperref}
\usepackage{xcolor}

\hypersetup{
  colorlinks=true,
  linkcolor=blue!50!black,
  citecolor=blue!50!black,
  urlcolor=blue!50!black
}

\newtheorem{theorem}{Theorem}[section]
\newtheorem{proposition}[theorem]{Proposition}
\newtheorem{lemma}[theorem]{Lemma}
\newtheorem{conjecture}[theorem]{Conjecture}
\newtheorem{problem}[theorem]{Problem}
\newtheorem{corollary}[theorem]{Corollary}

\theoremstyle{definition}
\newtheorem{definition}[theorem]{Definition}
\newtheorem{remark}[theorem]{Remark}
\newtheorem{example}[theorem]{Example}
\newtheorem{status}[theorem]{Status}

\newcommand{\graphon}{\mathcal W}
\newcommand{\E}{\mathbb E}
\newcommand{\1}{\mathbf 1}
\newcommand{\R}{\mathbb R}
\newcommand{\Z}{\mathbb Z}
\newcommand{\dd}{\,d}
\newcommand{\tk}{T_k}
\newcommand{\LS}{\mathrm{LS}}
\newcommand{\dist}{\operatorname{dist}}
\newcommand{\Hom}{\operatorname{Hom}}
\newcommand{\ord}{\operatorname{ord}}

\title[Progress on two graphon optimisation problems]{Progress on two graphon optimisation problems: counterexamples, corrected conjectures, and reduced variational targets}
\author{Working research note}
\date{\today}

\begin{document}

\begin{abstract}
This note records the current state of two graphon optimisation problems.  The first asks whether, for a particular six-vertex graph $H$, the minimiser of $t(H,W)$ at edge density $1-1/k$ is never the balanced $k$-partite graphon nor the constant graphon, and asks for an exact solution at $k=5$.  We verify the balanced and constant benchmark values, prove the proposed $k=5$ circular construction has density $24349/187500$, and give explicit competitors beating both benchmark graphons for $k=3,4,5$.  We also explain why the all-$k$ statement is not accessible by a local perturbation: for $k\ge 6$ the balanced $k$-partite graphon is first-order locally stable against all edge-density-preserving feasible perturbations.

The second problem asks for a convex/concave classification of the extremal curves
\[
   f_F(x)=\min\{t(F,W):t(K_2,W)=x\}.
\]
We show that the proposed extension to all larger graphs is false: $F=K_3\sqcup K_3$ already gives a curve that is neither convex nor concave on the interval $[1/2,2/3]$.  We formulate a repaired ``odd-prime'' direction that removes this product obstruction.  The clique subfamily is solved by Reiher's clique-density theorem and is piecewise concave.  However, two non-clique cases listed as concave in the original problem, namely $K_4$ with a pendant edge and $K_3\cup_{K_2}K_4$, do not have the usual Lovasz--Simonovits multipartite graphon as their minimiser.  We give explicit better graphons at edge density $17/25$, and then derive exact finite-dimensional variational formulae for the natural tripartite Turan-filling family.  The remaining main target is to prove, or refute, the global reduction to these filling families by flag algebras or symmetrisation.
\end{abstract}

\maketitle

\tableofcontents

\section{Notation and basic conventions}

A \emph{graphon} is a symmetric measurable function
\[
   W:[0,1]^2\to[0,1].
\]
For a finite graph $F=(V(F),E(F))$, its homomorphism density in $W$ is
\[
   t(F,W)=\int_{[0,1]^{V(F)}}\prod_{ij\in E(F)} W(x_i,x_j)\prod_{i\in V(F)}\dd x_i.
\]
In particular,
\[
   t(K_2,W)=\int_{[0,1]^2}W(x,y)\dd x\dd y
\]
is the ordered-pair edge density.  For a fixed graph $F$, define
\[
   f_F(x)=\min\{t(F,W):t(K_2,W)=x\},\qquad x\in[0,1].
\]
Existence of a minimiser follows from standard graphon compactness and lower semicontinuity of homomorphism densities.

The balanced complete $k$-partite graphon will be denoted by $T_k$.  It is obtained by partitioning $[0,1]$ into $k$ equal parts and putting value $0$ on each diagonal part and value $1$ between distinct parts.  Thus
\[
   t(K_2,T_k)=1-\frac1k.
\]
The constant graphon of edge density $p$ is denoted by $Q_p$; thus $Q_p(x,y)=p$ and
\[
   t(F,Q_p)=p^{e(F)}.
\]

We write
\[
   I_k=\left[1-\frac1k,\;1-\frac1{k+1}\right]
\]
for the usual Turan intervals.

\section{Problem 1: the six-vertex graph}

Let $H$ be the graph with vertex set $\{1,2,3,4,5,6\}$ and edge set
\[
\begin{split}
E(H)=\{&13,14,23,24,35,36,45,46,56\}.
\end{split}
\]
Thus $H$ has six vertices and nine edges.  The first problem asks about
\[
   \min\{t(H,W):t(K_2,W)=1-1/k\}.
\]

\subsection{The balanced $k$-partite and constant values}

\begin{lemma}[Balanced $k$-partite value]
For every integer $k\ge 1$,
\[
   t(H,T_k)=\frac{k(k-1)(k-2)(k^3-6k^2+14k-11)}{k^6}.
\]
\end{lemma}

\begin{proof}
A homomorphism from $H$ to $T_k$ is exactly a proper $k$-colouring of $H$.  First choose colours of vertices $5$ and $6$.  They must be distinct, giving $k(k-1)$ choices.  The vertices $3$ and $4$ must avoid both colours used on $5$ and $6$.

There are two cases.  If $3$ and $4$ receive the same colour, there are $k-2$ choices for this colour; then vertices $1$ and $2$ may each use any colour except that colour, giving $(k-1)^2$ choices.  If $3$ and $4$ receive distinct colours, there are $(k-2)(k-3)$ choices; then vertices $1$ and $2$ may each use any colour except the two colours used on $3,4$, giving $(k-2)^2$ choices.  Hence
\[
\begin{split}
\chi_H(k)&=k(k-1)\bigl((k-2)(k-1)^2+(k-2)(k-3)(k-2)^2\bigr)\\
&=k(k-1)(k-2)(k^3-6k^2+14k-11).
\end{split}
\]
Dividing by $k^6$ gives the stated homomorphism density.
\end{proof}

\begin{lemma}[Constant value]
For $Q_{1-1/k}\equiv 1-1/k$,
\[
   t(H,Q_{1-1/k})=\left(1-\frac1k\right)^9.
\]
\end{lemma}

\begin{proof}
The graph $H$ has nine edges, and every edge contributes the constant factor $1-1/k$.
\end{proof}

For reference, the first few benchmark values are
\[
\begin{array}{c|c|c}
 k & t(H,T_k) & t(H,Q_{1-1/k})\\ \hline
 3 & 8/243\approx 0.0329218 & (2/3)^9\approx 0.0260123\\
 4 & 39/512\approx 0.0761719 & (3/4)^9\approx 0.0750847\\
 5 & 2040/15625=0.13056 & (4/5)^9\approx 0.1342177.
\end{array}
\]

\subsection{The proposed circular construction at $k=5$}

Define
\[
   W_\circ(x,y)=\1\{0.1\le |x-y|\le 0.9\}.
\]
Equivalently, the zero set is the circular interval of radius $1/10$ around the diagonal.  Its measure is $1/5$, so
\[
   t(K_2,W_\circ)=\frac45.
\]

\begin{theorem}[Density of the circular construction]
The circular graphon $W_\circ$ satisfies
\[
   t(H,W_\circ)=\frac{24349}{187500}\approx 0.129861333\ldots .
\]
Consequently
\[
   t(H,W_\circ)<t(H,T_5)<t(H,Q_{4/5}).
\]
\end{theorem}

\begin{proof}
We approximate $W_\circ$ by finite cyclic step graphons.  Let $d$ be odd and set $N=5d$.  Partition $[0,1]$ into $N$ equal intervals indexed by $\Z/N\Z$.  Let $W_d$ be the step graphon with matrix
\[
   W_d(i,j)=1\quad\Longleftrightarrow\quad
   \dist_{C_N}(i,j)>\frac{d-1}{2}.
\]
For every row there are exactly $d$ zero entries, including the diagonal entry.  Hence
\[
   t(K_2,W_d)=1-\frac{d}{5d}=\frac45.
\]
Moreover $W_d\to W_\circ$ in $L^1$, and therefore $t(H,W_d)\to t(H,W_\circ)$.

A direct cyclic count gives
\[
   \#\Hom(H,W_d)=\frac{24349d^6+140d^4-9d^2}{12}.
\]
Here $\Hom(H,W_d)$ means the number of maps $\phi:V(H)\to\Z/N\Z$ such that all nine required pairs are $W_d$-edges.  Dividing by $N^6=(5d)^6$ gives
\[
   t(H,W_d)=\frac{24349+140d^{-2}-9d^{-4}}{187500}.
\]
Letting $d\to\infty$ gives
\[
   t(H,W_\circ)=\frac{24349}{187500}.
\]
Finally,
\[
   \frac{24349}{187500}<\frac{2040}{15625}<\left(\frac45\right)^9.
\]
Indeed $2040/15625=24480/187500$, and $24480>24349$.
\end{proof}

\begin{remark}[What this proves and what it does not prove]
The theorem proves that the balanced $5$-partite and constant graphons are not optimal at edge density $4/5$.  It does \emph{not} prove that $W_\circ$ is globally optimal.  The missing statement is the lower bound
\[
   t(H,W)\ge \frac{24349}{187500}\qquad
   \text{whenever }t(K_2,W)=4/5.
\]
This lower bound is the hard part of the $k=5$ problem and is a natural flag-algebra target.
\end{remark}

\subsection{Explicit competitors for $k=3,4,5$}

There is a simple one-parameter family that beats both the balanced and constant graphons for $k=3,4,5$.

Partition $[0,1]$ into $k$ equal parts and define
\[
   S_{k,a}(x,y)=
   \begin{cases}
      a, & x,y\text{ in the same part},\\[1mm]
      1-\dfrac{a}{k-1}, & x,y\text{ in different parts}.
   \end{cases}
\]
Then
\[
   t(K_2,S_{k,a})=1-\frac1k
\]
for every $a\in[0,1]$ for which $1-a/(k-1)\in[0,1]$.

For this family,
\[
   t(H,S_{k,a})=\frac1{k^6}\sum_{c_1,\ldots,c_6\in[k]}
        \prod_{ij\in E(H)}w(c_i,c_j),
\]
where
\[
   w(c_i,c_j)=
   \begin{cases}
      a, & c_i=c_j,\\
      1-a/(k-1), & c_i\ne c_j.
   \end{cases}
\]
Exact rational evaluation gives:
\[
\begin{array}{c|c|c|c}
 k & a & t(H,S_{k,a}) & \text{comparison}\\ \hline
 3 & 3/10 & \dfrac{792167976529}{31104000000000}\approx 0.02546836
   & <(2/3)^9<t(H,T_3)\\[2mm]
 4 & 1/5 & \dfrac{38420325073}{524880000000}\approx 0.07319830
   & <(3/4)^9<t(H,T_4)\\[2mm]
 5 & 1/20 & \dfrac{2186447178930623}{16777216000000000}\approx 0.13032241
   & <t(H,T_5)<(4/5)^9.
\end{array}
\]
Thus the statement ``neither balanced nor constant'' is rigorously verified for $k=3,4,5$.

\subsection{Why the all-$k$ statement needs a global argument}

The one-parameter perturbation above also explains why the problem changes at $k=6$.

\begin{proposition}[First-order local stability of $T_k$ for $k\ge 6$]
For the family $S_{k,a}$,
\[
   \left.\frac{d}{da}t(H,S_{k,a})\right|_{a=0}
   =\frac{6k^3-55k^2+142k-111}{k^5}.
\]
This derivative is negative for $k=3,4,5$ and positive for every $k\ge 6$.  Consequently, for $k\ge 6$, the balanced $k$-partite graphon is first-order locally stable against every edge-density-preserving feasible perturbation.
\end{proposition}

\begin{proof}
The displayed derivative follows by differentiating the finite sum for $t(H,S_{k,a})$ at $a=0$.  The polynomial
\[
   6k^3-55k^2+142k-111
\]
is negative at $k=3,4,5$ and positive at $k=6$.  Its derivative is
\[
   18k^2-110k+142,
\]
which is positive for $k\ge 5$, so the polynomial remains positive for all $k\ge 6$.

It remains to explain why this one-parameter derivative controls all first-order feasible directions at $T_k$.  By the symmetry of $T_k$, the first variation of $t(H,\cdot)$ at $T_k$ has one coefficient on within-part cells and one coefficient on cross-part cells.  A feasible one-sided perturbation can only add mass on within-part cells, where $T_k=0$, and remove mass from cross-part cells, where $T_k=1$.  The edge-density constraint says that the total added mass equals the total removed mass.  Hence every first-order feasible derivative is a constant multiple of the difference between the within-cell and cross-cell first-variation coefficients.  The family $S_{k,a}$ is exactly one such edge-density-preserving direction, so its derivative determines the sign for all first-order feasible directions.
\end{proof}

\begin{status}[Status of Problem 1]
The following points are solid.
\begin{enumerate}[label=(\roman*)]
\item The balanced and constant benchmark values are explicit.
\item At $k=5$, the circular graphon has the proposed value $24349/187500$ and beats both benchmarks.
\item For $k=3,4,5$, explicit rational competitors beat both benchmarks.
\item For $k\ge 6$, non-optimality of $T_k$ cannot be proved by an infinitesimal edge-density-preserving perturbation, because $T_k$ is first-order locally stable.
\end{enumerate}
The missing results are global lower and upper-bound statements.  The two most important targets are:
\begin{enumerate}[label=(\alph*)]
\item prove or disprove global optimality of the circular $k=5$ graphon;
\item find a genuinely global competitor, or a proof of optimality, for $k\ge 6$.
\end{enumerate}
\end{status}

\begin{remark}[Normalisation of circular thresholds]
To obtain edge density $1-1/k$ from a circular interval construction, the forbidden circular radius should be $1/(2k)$, not $1/k$.  Thus the literal analogue of the $k=5$ thresholds $0.1,0.9$ would use thresholds
\[
   \frac1{2k},\qquad 1-\frac1{2k}.
\]
For $k=3,4$, these would be $1/6,5/6$ and $1/8,7/8$, respectively.  This explains why replacing $0.1,0.9$ by $1/3,2/3$ or $1/4,3/4$ does not preserve the edge density.
\end{remark}

\section{Problem 2: the original extension is false}

The second problem asks for a convex/concave classification of $f_F(x)$ for graphs $F$ with no isolated vertices, no isolated edges, and at least one odd cycle.  It then asks to extend the classification to all graphs with more than five vertices.

The extension to all larger graphs is false.

\subsection{A counterexample: two disjoint triangles}

Let
\[
   F=K_3\sqcup K_3.
\]
This graph has six vertices, no isolated vertices, no isolated edges, and is non-bipartite.  For every graphon $W$,
\[
   t(F,W)=t(K_3,W)^2.
\]
Since $z\mapsto z^2$ is increasing on $[0,\infty)$,
\[
   f_F(x)=f_{K_3}(x)^2.
\]

On the Turan interval
\[
   I_2=\left[\frac12,\frac23\right],
\]
the triangle-density theorem gives the usual tripartite minimiser with part sizes
\[
   s,\qquad \frac{1-s}{2},\qquad \frac{1-s}{2},
   \qquad 0\le s\le \frac13.
\]
The edge and triangle densities are
\[
   x=\frac12+s-\frac32s^2
\]
and
\[
   f_{K_3}(x)=\frac32s(1-s)^2.
\]
Therefore
\[
   f_{K_3\sqcup K_3}(x)=\left(\frac32s(1-s)^2\right)^2.
\]
A direct calculation gives
\[
   \frac{d^2}{dx^2}f_{K_3\sqcup K_3}(x)
   =\frac{9(1-s)^2(4s-1)}{2(3s-1)}.
\]
For $0<s<1/4$, this is positive, while for $1/4<s<1/3$, it is negative.  Hence $f_{K_3\sqcup K_3}$ is neither convex nor concave on $I_2$.

\subsection{A family of product obstructions}

More generally, let
\[
   F_m=\underbrace{K_3\sqcup\cdots\sqcup K_3}_{m\text{ copies}},\qquad m\ge 2.
\]
Then
\[
   f_{F_m}(x)=f_{K_3}(x)^m.
\]
Using the same parameter $s$ on $I_2$,
\[
   f_{F_m}(x)=\left(\frac32s(1-s)^2\right)^m.
\]
A calculation gives
\[
   \frac{d^2}{dx^2}f_{F_m}(x)
   =
   \frac{m\left(\frac32s(1-s)^2\right)^m
     \bigl((3m-2)s-(m-1)\bigr)}
        {s^2(1-s)^2(3s-1)}.
\]
The sign changes at
\[
   s_0=\frac{m-1}{3m-2}\in\left(0,\frac13\right).
\]
Thus all $F_m$ with $m\ge 2$ have mixed curvature on a single Turan interval.

\begin{status}[Consequence]
Any correct extension of Problem 2 must exclude graphs with multiple independent non-bipartite pieces.  Merely requiring connectedness is probably not the right repair, because two odd blocks joined by a bridge should still exhibit a coupled product mechanism.
\end{status}

\section{A repaired structural direction: odd-prime graphs}

The counterexample suggests that the obstruction is not simply disconnectedness but the presence of multiple odd pieces.

\begin{definition}[Odd block and odd-prime graph]
A block of a graph is a maximal $2$-connected subgraph, with bridges treated as $K_2$-blocks.  A block is \emph{odd} if it contains an odd cycle.  Let $b_{\rm odd}(F)$ be the number of odd blocks of $F$.  A graph $F$ is called \emph{odd-prime} if it is connected, non-bipartite, and
\[
   b_{\rm odd}(F)=1.
\]
\end{definition}

The first repair of Problem 2 is therefore:

\begin{conjecture}[Odd-prime reduction principle]
Graphs with $b_{\rm odd}(F)\ge 2$ should not be expected to obey a pure convex/concave dichotomy.  For such graphs, the correct theory should decompose $f_F$ into products or coupled products of the extremal curves associated with the odd blocks.

For odd-prime graphs, a convex/concave classification may still exist, but the extremal templates need not be the same as the clique-density templates.  The right task is to identify, for each odd-prime family, the correct Turan-type filling construction and then prove the corresponding lower bound.
\end{conjecture}

This conjecture is deliberately weaker than the original all-graphs extension.  It says that the clique cases are only one part of the story.  Non-clique odd-prime graphs may have different extremal mechanisms.

\section{Clique cases: a solved concave subfamily}

The clique cases are completely understood by the clique-density theorem of Reiher \cite{ReiherCliqueDensity}.  We record the graphon formulation and the curvature calculation.

Fix $r\ge 3$.  On
\[
   I_k=\left[1-\frac1k,1-\frac1{k+1}\right],
\]
the Lovasz--Simonovits graphon has $k+1$ parts: one part of size $s$ and $k$ equal parts of size $(1-s)/k$, where
\[
   0\le s\le \frac1{k+1}.
\]
Its edge density is
\[
   x_k(s)=1-s^2-\frac{(1-s)^2}{k}.
\]
Its $K_r$-density is
\[
   g_{r,k}(s)=
   \frac{(k)_{r-1}}{k^{r-1}}(1-s)^{r-1}
   \left(
      \frac{k-r+1}{k}+\frac{(r-1)(k+1)}{k}s
   \right),
\]
where
\[
   (k)_{r-1}=k(k-1)\cdots(k-r+2).
\]
If $k+1<r$, then $g_{r,k}=0$.

\begin{theorem}[Clique curves are piecewise concave]
For every $r\ge 3$, $f_{K_r}$ is concave on every Turan interval $I_k$.
\end{theorem}

\begin{proof}
The clique-density theorem identifies $f_{K_r}$ with the Lovasz--Simonovits curve $g_{r,k}(s)$ on $I_k$.  We compute curvature in the parameter $s$.

First,
\[
   x_k'(s)=\frac{2}{k}\bigl(1-(k+1)s\bigr)>0
\]
in the interior of the interval.  Therefore the sign of $d^2g_{r,k}/dx^2$ is the sign of
\[
   g_{r,k}''(s)x_k'(s)-g_{r,k}'(s)x_k''(s).
\]
A direct calculation gives
\[
\begin{split}
&g_{r,k}''(s)x_k'(s)-g_{r,k}'(s)x_k''(s)\\
&\qquad =-
\frac{2r(r-1)(r-2)}{k^2}
\frac{(k)_{r-1}}{k^{r-1}}
(1-s)^{r-3}\bigl(1-(k+1)s\bigr)^2.
\end{split}
\]
This is non-positive.  Hence $f_{K_r}$ is concave on $I_k$.
\end{proof}

\begin{remark}
The clique case should be viewed as the solved model problem.  It gives exact curves, exact extremisers, and the desired concavity.  The non-clique cases below show that one cannot simply attach rooted bipartite decorations to clique extremisers and expect the same templates to remain optimal.
\end{remark}

\section{Two non-clique cases: the usual clique template is not optimal}

Let
\[
   P=K_4^\dagger
\]
denote $K_4$ with one pendant edge attached to one vertex of the $K_4$.  Let
\[
   J=K_3\cup_{K_2}K_4
\]
denote the graph obtained by gluing one edge of a triangle to one edge of a $K_4$.

Both graphs appear in the original list of proposed concave cases.  The following calculation shows that the usual Lovasz--Simonovits $4$-partite clique template is not optimal for either graph at edge density $17/25$.

\subsection{The Lovasz--Simonovits value at $x=17/25$}

On $I_3=[2/3,3/4]$, the clique-density template is the complete $4$-partite graphon with part sizes
\[
   s,\qquad \frac{1-s}{3},\qquad \frac{1-s}{3},\qquad \frac{1-s}{3}.
\]
Its edge density is
\[
   x(s)=\frac23+\frac{2s}{3}-\frac{4s^2}{3}.
\]
Setting $x=17/25$ gives
\[
   s=\frac{5-\sqrt{21}}{20}.
\]
For this graphon,
\[
   t(P)=\frac23s(1-s)^3
   =\frac{291}{10000}-\frac{7\sqrt{21}}{2000}
   \approx 0.013060985,
\]
and
\[
   t(J)=\frac{97}{5000}-\frac{7\sqrt{21}}{3000}
   \approx 0.008707323.
\]

\subsection{A better tripartite filling graphon}

Partition $[0,1]$ into three parts
\[
   A\sqcup B\sqcup C
\]
with
\[
   |A|=\alpha=\frac7{20},
   \qquad |B|=|C|=\frac{13}{40}.
\]
Put all cross-edges between distinct parts, put no edges inside $B$ or $C$, and inside $A$ put a $q$-regular triangle-free graphon with
\[
   q=\frac{11}{98}.
\]
For example, split $A$ into two equal subparts and put value $2q=11/49$ between the subparts and value $0$ inside each subpart.

The edge density is
\[
   t(K_2,W)=\frac12+\alpha-\frac32\alpha^2+\alpha^2q.
\]
Substituting $\alpha=7/20$ and $q=11/98$ gives
\[
   t(K_2,W)=\frac{533}{800}+\frac{11}{800}=\frac{17}{25}.
\]
For $P$, a direct rooted count gives
\[
   t(P,W)=\frac32\alpha^2(1-\alpha)^2q
          \left(\frac{3-\alpha}{2}+\alpha q\right).
\]
Thus
\[
   t(P,W)=\frac{1065207}{89600000}\approx 0.011888471,
\]
which is strictly smaller than the Lovasz--Simonovits value
\[
   \frac{291}{10000}-\frac{7\sqrt{21}}{2000}\approx 0.013060985.
\]
For $J$, the same graphon gives
\[
   t(J,W)=\alpha^2(1-\alpha)^2q
          \left(\frac32-\alpha+2\alpha q\right)
       =\frac{79937}{11200000}\approx 0.007137232,
\]
which is strictly smaller than
\[
   \frac{97}{5000}-\frac{7\sqrt{21}}{3000}\approx 0.008707323.
\]

\begin{status}[Consequence]
The two non-clique cases $P$ and $J$ cannot be solved by simply importing the clique-density extremiser.  Their extremisers, if the natural candidates are correct, involve a complete tripartite skeleton with a triangle-free filling inside one part.
\end{status}

\section{The tripartite Turan-filling reduction}

We now derive the exact variational formula inside the natural family suggested by the previous examples.  This part is rigorous as a reduction to a finite-dimensional optimisation problem.  The remaining global problem is to prove that every graphon minimiser can be chosen from this family.

\subsection{The family}

Let
\[
   [0,1]=A\sqcup B\sqcup C
\]
with
\[
   |A|=\alpha,
   \qquad |B|=|C|=\beta=\frac{1-\alpha}{2}.
\]
Put all cross-edges between $A,B,C$, put no edges inside $B$ or $C$, and put a triangle-free graphon $F$ inside $A$.  Normalize $F$ on a probability space of measure $1$ and write
\[
   q=t(K_2,F),
   \qquad d_F(u)=\int F(u,v)\dd v,
   \qquad D_2(F)=\int d_F(u)^2\dd u.
\]
By the graphon form of Mantel's theorem,
\[
   0\le q\le \frac12.
\]
The edge density of the full graphon is
\[
   x=\frac12+\alpha-\frac32\alpha^2+\alpha^2q.
\]
For fixed $x$, this gives
\[
   q=q_x(\alpha)=\frac{x-\frac12-\alpha+\frac32\alpha^2}{\alpha^2}.
\]
The constraint $0\le q_x(\alpha)\le 1/2$ is, for $x\in[2/3,3/4]$, equivalent to
\[
   \alpha\in A_x:=\left[
      \frac{1-\sqrt{3-4x}}2,
      \frac{1+\sqrt{3-4x}}2
   \right].
\]

\subsection{Counting $P=K_4^\dagger$}

\begin{lemma}[Density of $P$ in the filling family]
For the tripartite filling graphon above,
\[
   t(P,W)=\frac32\alpha^2(1-\alpha)^2
          \left[
             \frac{3-\alpha}{2}q+\alpha D_2(F)
          \right].
\]
Consequently, for fixed $\alpha$ and $q$, the minimum over triangle-free fillings $F$ is attained by a $q$-regular filling and equals
\[
   \frac32\alpha^2(1-\alpha)^2q
          \left(\frac{3-\alpha}{2}+\alpha q\right).
\]
\end{lemma}

\begin{proof}
Every $K_4$ in the full graphon must use two adjacent vertices inside $A$, one vertex in $B$, and one vertex in $C$.  The pendant edge may attach to the distinguished $K_4$ vertex.  If that distinguished vertex lies outside $A$, the pendant neighbour sees degree determined only by the complete cross-edges.  If it lies inside $A$, the contribution involves the $F$-degree of that vertex.  Summing over the four possible positions of the distinguished vertex gives the displayed formula.

Finally, Jensen's inequality gives
\[
   D_2(F)=\int d_F(u)^2\dd u\ge \left(\int d_F(u)\dd u\right)^2=q^2.
\]
Equality holds exactly for $q$-regular fillings.  Such triangle-free $q$-regular graphons exist for every $q\in[0,1/2]$: split the space into two equal parts and put value $2q$ across the two parts and $0$ inside them.
\end{proof}

Therefore the reduced variational problem for $P$ is
\[
\boxed{
   \Psi_P(x)=
   \min_{\alpha\in A_x}
   \frac32\alpha^2(1-\alpha)^2q_x(\alpha)
   \left(\frac{3-\alpha}{2}+\alpha q_x(\alpha)\right).
}
\]
After eliminating $q_x(\alpha)$,
\[
\boxed{
   \Psi_P(x)=
   \min_{\alpha\in A_x}
   \frac{3(1-\alpha)^2
     (2\alpha^2+\alpha+2x-1)
     (3\alpha^2-2\alpha+2x-1)}{8\alpha}.
}
\]

\subsection{Counting $J=K_3\cup_{K_2}K_4$}

\begin{lemma}[Density of $J$ in the filling family]
For the tripartite filling graphon above,
\[
   t(J,W)=\alpha^2(1-\alpha)^2
          \left[
             \left(\frac32-\alpha\right)q+2\alpha D_2(F)
          \right].
\]
Consequently, for fixed $\alpha$ and $q$, the minimum over triangle-free fillings $F$ is attained by a $q$-regular filling and equals
\[
   \alpha^2(1-\alpha)^2q
          \left(\frac32-\alpha+2\alpha q\right).
\]
\end{lemma}

\begin{proof}
Label the shared edge by $1,2$, the remaining triangle vertex by $3$, and the two remaining $K_4$ vertices by $4,5$.  The $K_4$ portion must again use two adjacent vertices in $A$, one vertex in $B$, and one vertex in $C$.  There are three cases according to whether the shared edge lies inside $A$, has exactly one endpoint in $A$, or lies outside $A$.  Summing these rooted contributions gives
\[
   \alpha^2(1-\alpha)^2
          \left[
             \left(\frac32-\alpha\right)q+2\alpha D_2(F)
          \right].
\]
The Jensen step and the construction of regular triangle-free fillings are the same as in the previous lemma.
\end{proof}

Thus the reduced variational problem for $J$ is
\[
\boxed{
   \Psi_J(x)=
   \min_{\alpha\in A_x}
   \alpha^2(1-\alpha)^2q_x(\alpha)
   \left(\frac32-\alpha+2\alpha q_x(\alpha)\right).
}
\]
After eliminating $q_x(\alpha)$,
\[
\boxed{
   \Psi_J(x)=
   \min_{\alpha\in A_x}
   \frac{(1-\alpha)^2
     (3\alpha^2-2\alpha+2x-1)
     (4\alpha^2-\alpha+4x-2)}{4\alpha}.
}
\]

\subsection{Stationary equations and algebraic branches}

For $P$, an interior stationary point satisfies
\[
\begin{split}
0={}&2\alpha^3q^2+9\alpha^3q-3\alpha^3
+2\alpha^2q^2-9\alpha^2q+13\alpha^2\\
&+4\alpha q-13\alpha+3,
\end{split}
\]
where
\[
   q=q_x(\alpha).
\]
Solving this quadratic equation for $q$ gives the relevant branch
\[
   q_P(\alpha)=
   \frac{-9\alpha^2+9\alpha-4+
   \sqrt{105\alpha^4-242\alpha^3+153\alpha^2+8\alpha-8}}
   {4\alpha(\alpha+1)}.
\]
This branch satisfies
\[
   q_P(1/3)=0,
   \qquad q_P(1/2)=1/2.
\]
It parametrises a candidate reduced optimum by
\[
   x_P(\alpha)=\frac12+\alpha-\frac32\alpha^2+
      \alpha^2q_P(\alpha),
   \qquad \alpha\in[1/3,1/2].
\]
The corresponding value is
\[
   \Psi_P(x_P(\alpha))=\frac32\alpha^2(1-\alpha)^2q_P(\alpha)
   \left(\frac{3-\alpha}{2}+\alpha q_P(\alpha)\right).
\]

For $J$, the corresponding stationarity equation is
\[
\begin{split}
0={}&4\alpha^3q^2+18\alpha^3q-6\alpha^3
+4\alpha^2q^2-24\alpha^2q+17\alpha^2\\
&+8\alpha q-14\alpha+3.
\end{split}
\]
The relevant branch is
\[
   q_J(\alpha)=
   \frac{-(3\alpha-2)^2+
   \sqrt{105\alpha^4-260\alpha^3+204\alpha^2-52\alpha+4}}
   {4\alpha(\alpha+1)}.
\]
Again
\[
   q_J(1/3)=0,
   \qquad q_J(1/2)=1/2.
\]
It parametrises the candidate reduced optimum by
\[
   x_J(\alpha)=\frac12+\alpha-\frac32\alpha^2+
      \alpha^2q_J(\alpha),
   \qquad \alpha\in[1/3,1/2],
\]
and
\[
   \Psi_J(x_J(\alpha))=
   \alpha^2(1-\alpha)^2q_J(\alpha)
   \left(\frac32-\alpha+2\alpha q_J(\alpha)\right).
\]
At the endpoints,
\[
   \Psi_P(2/3)=0,
   \qquad \Psi_P(3/4)=\frac9{128},
\]
and
\[
   \Psi_J(2/3)=0,
   \qquad \Psi_J(3/4)=\frac3{64}.
\]

\begin{remark}[Algebraic verification still to formalise]
The counting and Jensen reductions above are fully rigorous.  The last one-dimensional minimisation is elementary but should be formalised with a Sturm-sequence certificate before being quoted as a publication-ready theorem.  Numerically and symbolically, the branches $q_P$ and $q_J$ give the unique minimising branch on $[2/3,3/4]$.  A useful next step is to add the explicit Sturm certificates to this note or to a computer-verifiable appendix.
\end{remark}

\subsection{Numerical example at $x=17/25$}

Solving the reduced one-dimensional problems at
\[
   x=\frac{17}{25}
\]
gives
\[
   \alpha_P\approx 0.3489617543,
   \qquad q_P\approx 0.1125007820,
\]
and
\[
   \Psi_P(17/25)\approx 0.01188713928.
\]
For $J$,
\[
   \alpha_J\approx 0.3509647872,
   \qquad q_J\approx 0.1120315971,
\]
and
\[
   \Psi_J(17/25)\approx 0.007136528755.
\]
The hand-picked values $\alpha=7/20$ and $q=11/98$ used above are therefore very close to the reduced optima but not exactly optimal.

\section{Where the project currently stands}

\subsection{Solved or essentially settled pieces}

\begin{enumerate}[label=(\arabic*)]
\item For Problem 1, the balanced $k$-partite and constant values are explicit.
\item The proposed $k=5$ circular graphon has density $24349/187500$ and beats both benchmark graphons.
\item Explicit rational competitors beat both benchmark graphons for $k=3,4,5$.
\item For $k\ge 6$, the balanced $k$-partite graphon is first-order locally stable.  Thus any all-$k$ non-optimality proof must be global.
\item The original all-graphs extension in Problem 2 is false.  The counterexample $K_3\sqcup K_3$ has mixed curvature on $[1/2,2/3]$.
\item All clique curves $f_{K_r}$ are piecewise concave by Reiher's clique-density theorem.
\item The non-clique graphs $K_4^\dagger$ and $K_3\cup_{K_2}K_4$ are not minimised by the usual Lovasz--Simonovits $4$-partite clique template at edge density $17/25$.
\item Inside the natural tripartite Turan-filling family, the variational problems for $K_4^\dagger$ and $K_3\cup_{K_2}K_4$ reduce to explicit one-dimensional algebraic minimisations.
\end{enumerate}

\subsection{Main open questions}

\begin{problem}[Problem 1, $k=5$ lower bound]
Prove or disprove that
\[
   t(H,W)\ge \frac{24349}{187500}
   \qquad\text{whenever }t(K_2,W)=\frac45.
\]
If true, identify all minimising graphons up to weak isomorphism.  If false, find a better graphon.
\end{problem}

\begin{problem}[Problem 1, all $k$]
For every $k>2$, decide whether the minimiser at edge density $1-1/k$ is neither $T_k$ nor the constant graphon.  For $k\ge 6$, this cannot be answered by first-order perturbation around $T_k$; either a global competitor or a global optimality proof is required.
\end{problem}

\begin{problem}[Odd-block decomposition]
Develop a decomposition theory for $f_F$ when $F$ has multiple odd blocks.  The family $mK_3$ shows that products of odd pieces can destroy any simple convex/concave dichotomy.
\end{problem}

\begin{problem}[Global reduction for $P$ and $J$]
Prove or refute the following reduction on $[2/3,3/4]$:
\[
   f_P(x)=\Psi_P(x),
   \qquad f_J(x)=\Psi_J(x).
\]
Equivalently, show that every minimiser can be replaced by a complete tripartite graphon with a regular triangle-free filling inside one part, without increasing the $P$- or $J$-density.
\end{problem}

\begin{problem}[Curvature after the corrected reduction]
Once the correct extremal curves are identified for $P$ and $J$, determine their actual curvature.  Preliminary numerical evidence for the reduced curves suggests that the original small-graph convex/concave classification may need revision even before considering graphs with more than five vertices.
\end{problem}

\section{Promising research directions}

\subsection{Flag-algebra certificates}

The most direct path to the missing lower bounds is a flag-algebra computation.  For Problem 1, the target certificate is
\[
   t(H,W)-\frac{24349}{187500}\ge 0
   \quad\text{under }t(K_2,W)=4/5.
\]
For the non-clique Problem 2 cases, the target certificates are the inequalities
\[
   t(P,W)\ge \Psi_P(x),
   \qquad t(J,W)\ge \Psi_J(x),
   \qquad x=t(K_2,W),
\]
at least for $x\in[2/3,3/4]$.

These are parameter-dependent inequalities.  One feasible approach is to choose several rational values of $x$, solve the SDP numerically, guess exact certificates, and then interpolate the parameter dependence.

\subsection{Symmetrisation and stability}

The tripartite filling construction suggests a symmetrisation theorem:

\begin{quote}
Among graphons of edge density $x\in[2/3,3/4]$, a minimiser for $P$ or $J$ can be chosen to have a complete tripartite skeleton, with all nontrivial structure inside one part.
\end{quote}

A proof might combine stability for the $K_4$ supersaturation problem with a Zykov-type symmetrisation that preserves or decreases the rooted decoration density.  This would be conceptually preferable to a purely computational flag-algebra certificate.

\subsection{Rooted density inequalities}

For clique-decorated graphs, unrooted clique density is not enough.  For instance,
\[
   t(K_4^\dagger,W)=\int \kappa_4(x) d_W(x)\dd x,
\]
where $\kappa_4(x)$ is the rooted density of $K_4$ at the pendant root and $d_W(x)$ is the degree.  The examples above exploit correlations between rooted clique density and degree.  A promising direction is therefore to develop sharp rooted clique-density inequalities under an edge-density constraint.

\subsection{Formal algebraic verification}

The reduced formulas for $\Psi_P$ and $\Psi_J$ are explicit enough to be certified by exact algebra.  A concrete next step is to produce a small companion file, in Sage, Mathematica, or Lean, that verifies:
\begin{enumerate}[label=(\roman*)]
\item the stationary equations;
\item uniqueness of the minimising branch by Sturm sequences;
\item monotonicity of $x_P(\alpha)$ and $x_J(\alpha)$ on $[1/3,1/2]$;
\item the curvature signs of the reduced curves.
\end{enumerate}
This would make the finite-dimensional part robust and publication-ready.

\subsection{Revisiting the small-graph classification}

Because $K_4^\dagger$ and $K_3\cup_{K_2}K_4$ do not use the clique template, the original classification for graphs on at most five vertices should be rechecked.  A sensible workflow is:
\begin{enumerate}[label=(\arabic*)]
\item enumerate all non-bipartite graphs on at most five vertices with no isolated vertices and no isolated $K_2$ components;
\item for each graph, generate candidate Turan-filling templates on the first relevant intervals;
\item solve the reduced finite-dimensional problems exactly;
\item use flag algebras to certify or refute the candidates.
\end{enumerate}

\section{A compact checklist for collaborators}

Here is the shortest useful checklist for continuing the project.

\begin{enumerate}[label=\textbf{Task \arabic*.}]
\item Try to certify the $k=5$ circular value for Problem 1 by flag algebras.  The target constant is $24349/187500$ at edge density $4/5$.
\item Search for global competitors to $T_k$ for Problem 1 when $k\ge 6$.  Local perturbations will not find them.
\item Treat $K_3\sqcup K_3$ as the canonical obstruction to the original Problem 2 extension.
\item Do not assume that clique-decorated graphs are minimised by clique-density templates.  The cases $K_4^\dagger$ and $K_3\cup_{K_2}K_4$ already disprove that strategy.
\item For $K_4^\dagger$ and $K_3\cup_{K_2}K_4$, focus on proving or disproving the tripartite Turan-filling reduction.
\item Add exact Sturm certificates for the one-dimensional reduced minimisations.
\end{enumerate}

\appendix

\section{Exact finite-sum checker for the family $S_{k,a}$}

The following formula is useful for independent verification:
\[
   t(H,S_{k,a})=\frac1{k^6}\sum_{c_1,\ldots,c_6\in[k]}
        \prod_{ij\in E(H)}
        \left(a\1\{c_i=c_j\}+\left(1-\frac{a}{k-1}\right)\1\{c_i\ne c_j\}\right).
\]
This finite sum is the source of the rational values in Section~2.3 and of the derivative
\[
   \left.\frac{d}{da}t(H,S_{k,a})\right|_{a=0}
   =\frac{6k^3-55k^2+142k-111}{k^5}.
\]

\section{Derivative formulae for the reduced $P$ and $J$ problems}

For $P$, after eliminating $q$, define
\[
   \Phi_P(\alpha,x)=
   \frac{3(1-\alpha)^2
     (2\alpha^2+\alpha+2x-1)
     (3\alpha^2-2\alpha+2x-1)}{8\alpha}.
\]
Then
\[
\frac{\partial \Phi_P}{\partial \alpha}
=\frac{3(\alpha-1)}{8\alpha^2}N_P(\alpha,x),
\]
where
\[
\begin{split}
N_P(\alpha,x)={}&30\alpha^5-22\alpha^4+30\alpha^3x-19\alpha^3
-14\alpha^2x+9\alpha^2\\
&+4\alpha x^2-4\alpha x+\alpha+4x^2-4x+1.
\end{split}
\]
Substituting
\[
   x=\frac12+\alpha-\frac32\alpha^2+\alpha^2q
\]
into $N_P$ gives the stationarity equation in Section~7.4.

For $J$, define
\[
   \Phi_J(\alpha,x)=
   \frac{(1-\alpha)^2
     (3\alpha^2-2\alpha+2x-1)
     (4\alpha^2-\alpha+4x-2)}{4\alpha}.
\]
Then
\[
\frac{\partial \Phi_J}{\partial \alpha}
=\frac{(\alpha-1)}{2\alpha^2}N_J(\alpha,x),
\]
where
\[
\begin{split}
N_J(\alpha,x)={}&30\alpha^5-40\alpha^4+30\alpha^3x-\alpha^3
-20\alpha^2x+9\alpha^2\\
&+4\alpha x^2-4\alpha x+\alpha+4x^2-4x+1.
\end{split}
\]
Again substituting
\[
   x=\frac12+\alpha-\frac32\alpha^2+\alpha^2q
\]
gives the stationarity equation in Section~7.4.

\begin{thebibliography}{9}

\bibitem{ReiherCliqueDensity}
C. Reiher,
\emph{The clique density theorem},
Annals of Mathematics \textbf{184} (2016), 683--707.
Also available as arXiv:1212.2454.

\bibitem{BennettDudekLidickyPikhurkoC5}
P. Bennett, A. Dudek, B. Lidicky, and O. Pikhurko,
\emph{Minimizing the number of 5-cycles in graphs with given edge-density},
Combinatorics, Probability and Computing \textbf{29} (2020), 44--67.
Also available as arXiv:1803.00165.

\bibitem{LovaszBook}
L. Lovasz,
\emph{Large Networks and Graph Limits},
American Mathematical Society, 2012.

\end{thebibliography}

\end{document}
